\section{Preliminaries}
\label{sec:prelims}

% \subsection{A Review of Gaussian processes and \gpucb}
\subsection{Some Background Material}
\label{sec:prelims}

\textbf{Gaussian processes:}
A GP over a space $\Xcal$ is a random process describing functions from $\Xcal$ to $\RR$.
It is characterised by a mean function $\mu:\Xcal\rightarrow\RR$ and
a covariance kernel $\kernel:\Xcal^2\rightarrow\RR$.
If $\func\sim\GP(\mu,\kernel)$, then $f(x)$ is distributed normally
$\Ncal(\mu(x), \kernel(x,x))$ for all $x\in\Xcal$.
Suppose we are given $n$ observations from this GP, $\Dcal_n=\{(x_i, y_i)\}_{i=1}^n$ where
$x_i\in\Xcal$, $y_i = \func(x_i) + \epsilon_i \in\RR$ and $\epsilon_i\sim\Ncal(0,\eta^2)$.
Then the posterior process $\func|\Dcal_n$ is also a GP with some mean $\mu_n$ and
covariance $\kernel_n$, where
\begin{align*}
\hspace{-0.05in}
\mu_n(x) &= \kb^\top\Delta^{-1}Y, \hspace{0.35in}
\numberthis \label{eqn:gpPost} \\
\kernel_n(x,x') &= \kernel(x,x') - \kb^\top\Delta^{-1}\kb'.
\hspace{0.05in}
\end{align*}
Here, $Y\in\RR^n$ with $Y_i=y_i$,
$\kb,\kb'\in\RR^n$ with $\kb_i =
\kernel(x,x_i),\kb'_i=\kernel(x',x_i)$ and
$\Delta = \Kb + \eta^2I \in \RR^{n\times n}$ where
$\Kb_{i,j} = \kernel(x_i,x_j)$~\citep{rasmussen06gps}.
% We refer readers unfamiliar with GPs to chapter 2 of~\citet{rasmussen06gps} for
% the formulae for $\mu_n,\kernel_n$ and other properties of GPs.

\textbf{Radial kernels:}
Radial kernels are commonly used in the prior covariance of GPs;
% The squared exponential (SE) kernel and \matern{} kernels, which are popularly used in
% practice, are radial kernels.
some examples are the squared exponential (SE) and \matern{} kernels~\citep{rasmussen06gps}.
When using a radial kernel,
the prior covariance can be written as $\kernel(x,x') = \kernelscale\phi(\|x-x'\|)$ and
depends only on the distance between $x$ and $x'$.
Here, the scale parameter $\kernelscale$ 
captures the magnitude $\func$ could deviate from $\mu$.
The radial kernel, $\phi:\RR_{+}\rightarrow\RR_+$, is a decreasing function with
$\|\phi\|_\infty = \phi(0) = 1$.
In this paper, we will use the SE kernel in a running example to convey the
intuitions in our methods. 
For the SE kernel, $\phi(r) = \radialh(r) = \exp(-r^2/(2h^2))$, where $h\in\RR_+$ is
the bandwidth parameter and controls the smoothness of the GP.
% It is well known that 
When $h$ is large, the samples drawn from the GP tend to
be smoother as illustrated in Fig.~\ref{fig:gpBWSample}.
We will reference this observation frequently in the text.

\insertFigBWSample

\textbf{\gpucb:}
The Gaussian Process Upper Confidence Bound (\gpucb) algorithm
of~\citet{srinivas10gpbandits}  is a method for bandit optimisation,
which, like many other \bayos methods, models $\func$ as a sample from a Gaussian process.
At time $t$, the next point $\xt$ for evaluating $\func$ is chosen via the following
procedure.
First, we construct an upper confidence bound
$\utilt(x) = \mutmo(x) + \betath \sigmatmo(x)$ for the GP. 
$\mutmo$ is the posterior mean of the GP conditioned on the previous $t-1$ evaluations
and $\sigmatmo$ is the posterior standard deviation.
Then the next point is chosen by maximising $\utilt$, i.e.
$\xt=\argmax_{x\in\Xcal}\utilt(x)$.
The $\mutmo$ term encourages an \emph{exploitative} strategy -- in that we want to query
regions where we already believe  $\func$ is high -- and $\sigmatmo$ encourages
an \emph{exploratory} strategy -- in that we want to query where we are uncertain
about $\func$ so that we do not miss regions which have not been queried yet.
$\betat$, which is typically increasing with $t$, controls the trade-off between
exploration and exploitation.


\subsection{Problem Set Up}
\label{sec:prelims}

% In stochastic bandit optimisation, our goal is to maximise a function
Our goal in bandit optimisation is to maximise a function
$\func:\Xcal\rightarrow \RR$, over a domain $\Xcal$.
When we evaluate $\func$ at $x\in\Xcal$ we observe $y=\func(x) + \epsilon$ where
$\EE[\epsilon] = 0$.
Let $\xopt \in \argmax_{x\in\Xcal}\func(x)$ be a maximiser of $\func$ and
$\funcopt = \func(\xopt)$ be the maximum value.
% $\func$ is accessible only via noisy point evaluations and has no gradient information.
An algorithm for bandit optimisation is a sequence of points $\{\xt\}_{t\geq0}$, where
at time $t$, the algorithm chooses to evaluate $\func$ at $\xt$ based on previous queries
and observations $\{(\xtt{i},\ytt{i})\}_{i=1}^{t-1}$.
After $n$ queries to $\func$,
its goal is to achieve small \emph{simple regret} $\Sn$, as defined below.
\begin{align*}
\Sn = \min_{t=1,\dots,n} \funcopt - \func(\xt).
\numberthis
\label{eqn:SnDefn}
\end{align*}
\textbf{Multi-fidelity Bandit Optimisation:}
In this work, we will let $\func$ be a slice of a function $\gunc$ that lies
in a larger space.
In particular, we will assume the existence of a fidelity space $\Zcal$ and a function 
$\gunc:\Zcal\times\Xcal \rightarrow \RR$ defined on the product space of the fidelity
space and domain.
The function $\func$ which we wish to maximise is related to $\gunc$ via
$\func(\cdot) = \gunc(\zhf, \cdot)$, where $\zhf\in\Zcal$.
For instance, in the hyper-parameter tuning example from Section~\ref{sec:intro},
$\Zcal = [1,\Nmax]\times[1,\Tmax]$ and $\zhf = [\Nmax, \Tmax]$.
Our goal is to find the maximiser $\xopt = \argmax_x \func(x) = \argmax_x \gunc(\zhf, x)$.
We have illustrated this set up in Fig.~\ref{fig:fidelSpace}.
In the rest of the manuscript, the term ``fidelities'' will refer to points $z$
in the fidelity space $\Zcal$.
\insertFigFidelSpace

The multi-fidelity framework is attractive when the following two conditions are true
about the problem.
\vspace{-0.10in}
% \begin{itemize}[label=-\hspace{-0.02in},leftmargin=0.13in]
\begin{enumerate}[leftmargin=0.19in]
\item
\ul{There exist fidelities $z\in\Zcal$ where evaluating $\gunc$ is cheaper than
evaluating at $\zhf$.}
To this end, we will associate a \emph{known} cost function $\cost:\Zcal\rightarrow
\RR_+$.
In the hyper-parameter tuning example, $\cost(z) = \cost(N,T) = \bigO(N^2T)$.
It is helpful to think of $\zhf$ as being the most expensive fidelity,
i.e. maximiser of $\cost$, and that $\cost(z)$ decreases as we move away from
$\zhf$. 
However, this notion is strictly not necessary for our algorithm
or results.\hspace{-0.1in}
\vspace{-0.06in}
\item
\ul{The cheap $\gunc(z,\cdot)$ evaluation gives us information about
$g(\zhf,\cdot)$.}
This is true if $\gunc$ is smooth across the fidelity space
as illustrated in Fig.~\ref{fig:fidelSpace}.
As we will describe shortly, this can be achieved by modelling $\gunc$ as a Gaussian process
with an appropriate kernel for $\Zcal$.
% \end{itemize}
\end{enumerate}
\vspace{-0.12in}

In the above set up, a multi-fidelity algorithm is a sequence of query-fidelity pairs
$\{(\zt,\xt)\}_{t\geq 0}$
where, at time $t$, the algorithm chooses $\zt\in\Zcal$ and $\xt\in\Xcal$ based on
previous fidelity-query-observation triples $\{(\ztt{i},\xtt{i},\ytt{i})\}_{i=1}^{t-1}$.
Here, $\ytt{i} = g(\ztt{i},\xtt{i}) + \epsilon$ where $\EE[\epsilon] = 0$.

\textbf{Multi-fidelity Simple Regret:}
We provide bounds on the simple regret $S(\capital)$ of a multi-fidelity optimisation method
after it has spent capital $\capital$ of a resource.
Following~\citet{srinivas10gpbandits,kandasamy16mfbo} we
will aim to provide \emph{any capital} bounds, meaning that an algorithm would be expected
to do well for \emph{all} values of (sufficiently large) $\capital$.
Say we have made $N$ queries to $g$ within capital $\capital$,
i.e. $N$ is the \emph{random} quantity such that
$N = \max\{n\geq 1: \sum_{t=1}^n \cost(\zt) \leq \capital\}$.
While the cheap evaluations at $z\neq \zhf$ are useful in guiding search 
for the optimum of $\gunc(\zhf,\cdot)$, there is no reward for optimising a cheaper
$g(z,\cdot)$. Accordingly we define the simple regret after capital $\capital$ as,
\vspace{-0.02in}
\begin{align*}
S(\capital) \,=
\begin{cases}
\displaystyle\min_{\substack{t=\,1,\dots,N\\t\,:\,\zt = \zhf}}\;
  \funcopt - \func(\xt) % \hspace{0.01in}
    &\text{if we have queried at $\zhf$,} \\
+\,\infty & \text{otherwise.}
\end{cases}
\\[-0.15in]
% \numberthis
% \label{eqn:SLambdaDefn}
\end{align*}
\vspace{-0.35in}

This definition reduces to the single fidelity definition~\eqref{eqn:SnDefn} when we only
query $g$ at $\zhf$.
It is also similar to the definition in~\citet{kandasamy16mfbo}, but unlike them, we do
not impose additional boundedness constraints on $\func$ or $\gunc$.

Before we proceed, we note that it is customary in the bandit literature to analyse
\emph{cumulative regret}. However, the definition of cumulative regret depends
on the
application at hand~\citep{kandasamy16mfbandit} and the results in this paper can be
extended to to many sensible notions of cumulative regret.
However, both to simplify exposition and since our focus in this paper is optimisation,
we stick to simple regret.

\textbf{Assumptions:}
As we will be primarily focusing on continuous and compact domains and fidelity spaces,
going forward we will assume, without any loss of generality, that $\Xcal = [0, 1]^d$ and
$\Zcal = [0, 1]^p$.
We discuss non-continuous settings briefly at the end of Section~\ref{sec:mfgpucb}.
In keeping with similar work in the Bayesian optimisation literature, we will
assume
$\gunc\sim\GP(\zero,\kernel)$ and upon querying at $(z,x)$ we observe
$y = g(z,x) + \epsilon$ where $\epsilon\sim\Ncal(0,\eta^2)$.
$\kernel:(\Zcal\times\Xcal)^2\rightarrow \RR$ is
the prior covariance defined on the product space.
In this work, we will study exclusively $\kernel$ of the following form,
\begin{align*}
\hspace{-0.1in}
\kernel([z,x],[z',x']) \,=\,
\kernelscale\,\phiz(\|z-z'\|)\, \phix(\|x-x'\|).
\label{eqn:mfkernel}
\numberthis
\end{align*}
Here, $\kernelscale\in\RR_+$ is the scale parameter and $\phiz,\phix$ are radial kernels
defined on $\Zcal,\Xcal$ respectively.
The fidelity space kernel $\phiz$ is an important component in this work.
It controls the
smoothness of $\gunc$ across the fidelity space and hence determines how much
information the lower fidelities provide about $g(\zhf,\cdot)$.
For example, suppose that $\phiz$ was a SE kernel.
A favourable setting for
a multi-fidelity method would be for $\phiz$ to have a large bandwidth $\hz$
as that would imply that $g$ is very smooth across $\Zcal$.
We will see that $\hz$ determines the behaviour and theoretical guarantees
of \mfgpucbtwos in a natural way when $\phiz$ is the SE kernel.
To formalise this, we will  define the following function
$\inffunc:\Zcal\rightarrow[0,1]$.
\begin{align*}
\inffunc(z) = \sqrt{1 - \phiz(\|z-\zhf\|)^2}
\label{eqn:inffunc}
\numberthis
\end{align*}
One interpretation of $\inffunc(z)$ is that it measures the \emph{gap} in
information about $g(\zhf,\cdot)$ when we query at $z\neq\zhf$.
That is, it is the price we have to pay, in information, for querying at a cheap fidelity.
Observe that $\inffunc$ increases when we move away from $\zhf$ in the fidelity space.
For the SE kernel, it can be shown $\inffunc(z) \approx \frac{\|z-\zhf\|}{\hz}$.
For large $\hz$, $\gunc$ is smoother across $\Zcal$ and we can expect the lower fidelities
to be more informative about $\func$;  as expected $\inffunc$ is small for large $\hz$.
If $\hz$ is small and $\gunc$ is not smooth, the loss of information $\inffunc$ is large
and lower fidelities are not as informative.

% \todo{@samy, see comment in the code below}
%%% @samy: when defining functions, dont forcibly include space. The right way to do it is to define the command as usual and call the function with trailing parentheses. This is because Latex, by default, assumes that any character following \bayos is an argument to the function \bayos, by default. This is why it messes up spaces sometimes. So, the right way to do it is to always call functions as '\bayos{}' (I have done this in this case).}

Before we present our algorithm for the above setup, we will introduce notation for the posterior GPs for $\gunc$ and $\func$.
Let $\Dcal_n = \{(\ztt{i},\xtt{i},\ytt{i})\}_{t=1}^n$ be $n$ fidelity, query,
observation values from the GP $\gunc$, where $\ytt{i}$ was observed when evaluating
$g(\ztt{i},\xtt{i})$.
We will denote the posterior mean and standard
deviation of $\gunc$ conditioned on $\Dcal_n$ by $\nutt{n}$ and $\tautt{n}$ respectively
($\nutt{n},\tautt{n}$ can be computed from~\eqref{eqn:gpPost} by replacing
$x\leftarrow[z,x]$).
Therefore $\gunc(z,x)|\Dcal_n\sim\Ncal(\nutt{n}(z,x), \tausqtt{n}(z,x))$ for all
$(z,x)\in\Zcal\times\Xcal$.
We will further denote 
\begin{align*}
\mutt{n}(\cdot) = \nutt{n}(\zhf,\cdot), \hspace{0.2in}
\sigmatt{n}(\cdot) = \tautt{n}(\zhf,\cdot),
\label{eqn:munsigman}
\numberthis
\end{align*}
to be the posterior mean and standard
deviation of $\gunc(\zhf,\cdot) = \func(\cdot)$.
It follows that $\func|\Dcal_n$ is also a GP and satisfies
$\func(x)|\Dcal_n \sim \Ncal(\mutt{n}(x), \sigmasqtt{n}(x))$ for all $x\in\Xcal$.

